\documentclass[11pt,letterpaper]{article}

% ============================================================================
% The Trajectory of Code — self-contained preamble (no external .sty needed)
% ============================================================================
\usepackage[margin=1in]{geometry}
\usepackage{amsmath,amssymb}
\usepackage{booktabs}
\usepackage{enumitem}
\usepackage{xcolor}
\usepackage[hidelinks]{hyperref}
\usepackage{microtype}

\setlist{itemsep=2pt,topsep=3pt}

\newcommand{\doctitle}[1]{{\LARGE\bfseries #1\par}\vspace{0.6em}}
\newcommand{\survival}{\ensuremath{S}}
\newcommand{\shape}{\ensuremath{\sigma}}

\begin{document}

% ---- Cover ------------------------------------------------------------------
\begin{titlepage}
\vspace*{2cm}
\begin{center}
\doctitle{The Trajectory of Code}
{\Large Survival-Weighted Learning from Version-Control History\\for Agentic Software Models}\par
\vspace{2cm}
{\large Hanzo AI}\par
\vspace{0.3em}
{Hanzo Industries Inc \quad(Techstars '17)}\par
{\small research@hanzo.ai \quad ORCID: --- \quad @zeekay}\par
\vspace{2cm}
{\small HAI-2026-014 \quad\textbullet\quad Version 0.1 (Draft) \quad\textbullet\quad \textsc{Internal / Private}}\par
\vspace{1cm}
\end{center}
\begin{abstract}
\noindent We argue that the most valuable training signal for agentic software
models is not the agent's transcript nor the final source tree, but the
\emph{longitudinal history of the code itself}. A version-control graph is a
record of natural selection: every line, hunk, and structural \emph{shape} that
an author commits is subsequently either preserved, rewritten, or deleted by the
future. Survival therefore constitutes a time-validated, annotation-free quality
label at massive scale. We formalize \emph{code survival} $\survival$ (the
duration a unit of code persists before it is changed or removed), extend it from
lines to abstract structural shapes ($\shape$), and show that survival-weighted
objectives let a model learn to generate code that \emph{lasts} rather than code
that merely parses or passes a single test. We describe a decomplected extraction
pipeline that fuses three co-registered layers---agent \emph{trajectories}
($\sim$5{,}190 Claude Code sessions), verified \emph{diffs} (file-history and git),
and \emph{outcomes} (survival, test status, and a per-commit \textsc{human}/\textsc{ai}
label)---across a heterogeneous fleet of developer machines, with mandatory secret
redaction and provenance versioning. The result is a corpus of the form
\textit{(instruction $\to$ agent trajectory $\to$ verified diff $\to$ survival
outcome)}, which supports supervised, preference, and outcome-verified
reinforcement objectives, and a novel evaluation regime---\emph{survival
holdout}---that scores a model on whether its generated changes would have endured.
\end{abstract}
\end{titlepage}

\tableofcontents
\newpage

% ============================================================================
\section{Thesis: time is the annotator}
% ============================================================================

Public code datasets are overwhelmingly \emph{cross-sectional}: a snapshot of a
repository at one instant, scraped and tokenized. Agentic benchmarks add the
task dimension (an instruction plus a repository state), and instruction-tuned
corpora add human preference labels. All three ask the same implicit question:
\emph{what code was written?} None asks the question that a working engineer
actually cares about: \emph{what code should have been written---and how do we
know?}

A version-control history answers that question for free, because it contains
the future. Every unit of code an author commits is, over the following weeks
and years, subjected to a merciless empirical test: it is kept, edited, or
deleted by real subsequent work. Code that survives dozens of refactors was
\emph{validated by time}; code rewritten the next morning is a labeled negative.
No annotator, rubric, or reward model produced these labels---the history did,
and it did so at the scale of every commit ever made.

This paper develops that observation into a training methodology. Our central
claims are:

\begin{enumerate}
\item \textbf{Survival is a quality signal.} The persistence of a unit of code
is a time-validated, annotation-free proxy for its quality, correctness, and
architectural fit.
\item \textbf{Shapes, not lines, are the unit of interest.} Line survival is
noisy; the durable signal lives at the level of \emph{shapes}---function
signatures, module boundaries, abstractions, and interfaces. Some shapes last
far longer than others, and those are the good bones.
\item \textbf{Trajectory beats snapshot.} A model should learn how code
\emph{moves}---the ordered sequence of states a file passes through---not merely
its terminal form.
\item \textbf{The three layers must be co-registered.} Agent reasoning (the
\textsc{how}), the verified diff (the \textsc{what}), and the survival outcome
(the \textsc{result}) are individually useful and jointly transformative: they
yield grounded tuples \textit{(instruction $\to$ trajectory $\to$ diff $\to$
outcome)} that no single-layer corpus can express.
\end{enumerate}

% ============================================================================
\section{Formalism: code survival}
% ============================================================================

\subsection{Line and hunk survival}

Let a repository's history be an ordered sequence of commits
$c_0, c_1, \dots, c_T$. For a unit of code $u$ (a line, a hunk, or a shape)
introduced at commit $c_b$ (\emph{birth}) and first modified or removed at
commit $c_d$ (\emph{death}), define its survival as the wall-clock duration

\[
\survival(u) \;=\; t(c_d) - t(c_b),
\qquad
E(u) \;=\; d - b,
\]

where $E(u)$ is the number of intervening commits the unit \emph{survived
through} (its edit-distance longevity). A unit still present at the head commit
$c_T$ is right-censored; we treat it as alive with a lower bound
$\survival(u) \ge t(c_T) - t(c_b)$ and handle censoring with standard
survival-analysis estimators (Kaplan--Meier for population curves; Cox models
when conditioning on features).

Line survival is recoverable from \texttt{git} via incremental blame reconstruction
(equivalently, the machinery behind code half-life analyses): replay the diff
stream, assign every added line a birth, and record its death when a later diff
deletes or replaces it. The output is a per-line record
$\langle \text{repo}, \text{path}, \text{birth}, \text{death}, \survival, E,
\text{author\_type}\rangle$.

\subsection{From lines to shapes}

Line survival is high-variance: reformatting, renaming, and mechanical edits
kill lines without killing \emph{ideas}. The load-bearing signal is the survival
of \emph{shapes}---structural units recovered from the abstract syntax tree:
function signatures, type/class definitions, module and package boundaries,
public interfaces, and recurring design patterns. A shape $\shape$ is identified
across commits by a stable structural key (e.g.\ a normalized signature or an
interface hash) rather than by text, so it can survive an internal rewrite that
replaces every one of its lines. We then measure $\survival(\shape)$ exactly as
above.

The empirical claim is that the distribution of $\survival(\shape)$ is heavy-tailed
and \emph{informative}: a small fraction of shapes---clean abstractions,
well-chosen interfaces---persist across the entire visible history, while the
majority churn. Training a generator to prefer the long-lived tail is training
it toward durable architecture, which is precisely the competence that
cross-sectional corpora fail to teach.

\subsection{Durability label}

Given a survival distribution we define a per-unit durability label by upper and
lower quantiles $q_{\text{hi}}, q_{\text{lo}}$ of $\survival$ (computed
within-repository and within-language to control for base rates):

\[
\ell(u) =
\begin{cases}
\textsc{durable} & \survival(u) \ge q_{\text{hi}}\\
\textsc{ephemeral} & \survival(u) \le q_{\text{lo}}\\
\textsc{neutral} & \text{otherwise.}
\end{cases}
\]

\textsc{ephemeral} units are not discarded---they are the negative class, and
they are as valuable as the positives.

% ============================================================================
\section{The three co-registered layers}
% ============================================================================

The corpus fuses three sources, each answering a different question, joined on a
common \emph{(repo, commit, timestamp)} key.

\begin{table}[h]
\centering
\begin{tabular}{@{}lll@{}}
\toprule
\textbf{Layer} & \textbf{Source} & \textbf{Teaches} \\
\midrule
\textsc{How} (reasoning + tools) & Claude Code session transcripts & the agent trajectory:\\
& ($\sim$5{,}190 sessions, $\sim$4.3\,GB) & task $\to$ think $\to$ tool call $\to$ observe \\[4pt]
\textsc{What} (verified change) & file-history + git diffs & the actual before$\to$after code \\[4pt]
\textsc{Result} (label + reward) & survival $+$ CI $+$ human/ai tag & did it last? did it pass? who wrote it? \\
\bottomrule
\end{tabular}
\caption{The three data layers. Individually useful; jointly, they yield grounded
\textit{(instruction $\to$ trajectory $\to$ diff $\to$ outcome)} tuples.}
\end{table}

The join is the value. Public agentic datasets have the \textsc{how} or the
\textsc{what}; almost none carry the \textsc{result} at the granularity of
survival, and none is additionally labeled by author type. Because our
transcripts and commits are co-located on the same machines, each agent
trajectory can be linked to the commit(s) it produced (by timestamp proximity,
repository, and the version-control tool calls embedded in the trajectory), and
each resulting hunk inherits a survival outcome. We can therefore ask, and
learn from, the question no other corpus can pose: \emph{which agent decisions
produce code that lasts?}

\subsection{The human/ai axis}

Every commit carries an author-type label---\textsc{human} (the developer's own
hands, including historical corporate identities) or \textsc{ai} (an autonomous
agent account). This axis is not decoration; it lets survival be conditioned on
authorship, quantifying (for example) whether AI-authored abstractions survive
at human rates, and isolating the trajectory features that predict durability
regardless of author. In the reference dataset the \textsc{ai} account first
appears in late 2024 and, within eighteen months, produces contributions at
several times the \textsc{human} rate---so the corpus also documents, in real
data, the transition of authorship from human to agent.

% ============================================================================
\section{Construction pipeline}
% ============================================================================

The pipeline is deliberately decomplected---one directed path, one artifact per
stage, one canonical implementation:

\[
\textbf{extract} \to \textbf{normalize} \to \textbf{link} \to \textbf{redact}
\to \textbf{format} \to \textbf{version/store}.
\]

\begin{description}[leftmargin=1.2em]
\item[Extract.] Two collectors. (i) A \emph{session} collector reads
\texttt{\textasciitilde/.claude/projects/*.jsonl} into structured trajectories.
(ii) A \emph{history} collector replays each repository's git graph into the
per-line/per-shape survival records of Section~2.
\item[Fleet aggregation.] Development is spread across a heterogeneous fleet of
machines, each holding a partial \texttt{\textasciitilde/.claude} and partial
working copies. The session collector runs \emph{per host} and ships its
sanitized shard to a central store; shards are deduplicated by session id and
content hash so re-runs are idempotent and no machine is a single point of
truth.
\item[Normalize.] One canonical schema:
\texttt{trajectory\{task, turns[\{think, tool\_calls, tool\_results\}], outcome\}}
and \texttt{unit\{repo, path, kind, birth, death, survival, edits, author\_type,
durable\}}.
\item[Link.] Join trajectory $\leftrightarrow$ commit $\leftrightarrow$ survival
on \emph{(repo, commit, time)}.
\item[Redact.] \textbf{Mandatory and non-negotiable.} Session transcripts and git
history are known to contain live credentials (private keys, cloud and service
tokens). Redaction runs before any artifact leaves its host: secret scanning,
key/token/\texttt{BEGIN}-block scrubbing, PII removal, and a private/public
license gate. Failed or aborted sessions are \emph{labeled}, not dropped.
\item[Format.] Emit SFT, preference, trajectory, and evaluation splits
(Section~5).
\item[Version/store.] All artifacts and metadata land in a single versioned
dataset (\texttt{zen-agentic-dataset}); large blobs are held out of git and
backed up to a hub (HuggingFace), with a private repository for the full corpus
and a public repository for the redacted, license-clean subset. The survival and
authorship metadata form the provenance layer for every record.
\end{description}

% ============================================================================
\section{Training and evaluation formulations}
% ============================================================================

Survival is not merely a filter; it is an objective.

\begin{description}[leftmargin=1.2em]
\item[Supervised (durable targets).] Standard next-token / next-action loss,
with examples weighted by $w(u) = f(\survival(u))$ so the model imitates
long-lived code more strongly than churn.
\item[Preference / DPO.] Pairs $(\textsc{durable}, \textsc{ephemeral})$ drawn
from the same file and neighborhood give a natural preference signal:
\emph{prefer the change that would have survived}.
\item[Outcome-verified RL.] The reward is the join of \emph{did it pass}
(CI/test status) and \emph{did it last} (survival), yielding a dense,
non-gameable signal for policies that act over tool-use trajectories.
\item[Trajectory modeling.] Train directly on state transitions
$s_t \to s_{t+1}$ (how a file evolves), teaching refactoring and growth rather
than one-shot synthesis.
\item[Survival-holdout evaluation.] A benchmark that scores a model on whether
its generated change matches (or would plausibly survive as well as) the change
that historically endured at that point in the file's trajectory. Unlike
single-test pass/fail, this rewards durable design, not merely green pipelines.
\end{description}

% ============================================================================
\section{Why survival-weighting produces better agents}
% ============================================================================

Imitation learning on cross-sectional code teaches a model to reproduce the
average of what was written---including the large fraction that was immediately
undone. Outcome-verified benchmarks improve on this but reward only the narrow
target of passing a fixed test, which is gameable and silent about architecture.
Survival-weighting adds the missing axis: it rewards code that continued to be
\emph{correct enough, clear enough, and well-factored enough that future work
built on it rather than around it}. That is the closest available analogue to a
natural-selection objective for software, and it is obtainable at the scale of a
career's version-control history rather than the scale of hand-written rubrics.
Because the same corpus carries the agent's reasoning and the author-type label,
it uniquely supports the question that matters for building coding agents:
\emph{which patterns of agent behavior yield durable code?}

% ============================================================================
\section{Privacy, provenance, and ethics}
% ============================================================================

The corpus is derived from private development history and therefore carries real
obligations. Secret redaction is a hard gate, enforced before egress and verified
in tests. Private-repository code is segregated under the private dataset with
license controls; only redacted, license-clean records enter the public split.
Every record retains provenance---source host, session id, commit, author type,
and pipeline version---so any datum can be traced, re-derived, or withdrawn. The
methodology stores \emph{all} exported data and metadata under one versioned home
with fleet-wide backup, so the corpus is durable, auditable, and reproducible.

% ============================================================================
\section{Conclusion}
% ============================================================================

The history of code is a longitudinal experiment that has already been run, at
enormous scale, with time as the annotator. Reading survival as a label---and
elevating the unit of interest from lines to the shapes that last---turns a
developer's version-control graph into a uniquely valuable training corpus for
agentic software models: one that teaches not what code was written, but what
code deserved to endure. Co-registered with agent trajectories and an
author-type axis, it lets us train and measure agents against the only judge
that ever mattered, which is the future.

\end{document}
